\chapter{PHÂN TÍCH YÊU CẦU VÀ NGHIỆP VỤ}
\section{Tổng quan nghiệp vụ}
\subsection{Nghiệp vụ trải nghiệm nước hoa}
% Khách xem mùi hương, chọn sản phẩm, xác nhận giá, thanh toán và nhận xịt.
\subsection{Nghiệp vụ quản lý vận hành}
% Quản lý thương hiệu, người dùng, máy, khe, tồn kho, cảnh báo, bảo trì.
\section{Các tác nhân và vai trò}
\subsection{Quản trị viên nền tảng (Platform Super Administrator)}
\subsection{Quản trị viên thương hiệu (Brand Administrator)}
\subsection{Quản lý vận hành (Operations Manager)}
\subsection{Kỹ thuật viên (Technician)}
\subsection{Nhân viên kho (Inventory Staff)}
\subsection{Khách hàng (Customer/Kiosk User)}
% Phân biệt vai trò trong đề xuất với vai trò thực tế của phiên bản triển khai.
\section{Yêu cầu chức năng}

\subsection{Quản lý thương hiệu, người dùng và phân quyền}
\begin{itemize}
  \item Quản trị viên nền tảng tạo, kích hoạt, tạm ngưng thương hiệu (tenant) và tạo tài khoản Quản trị viên thương hiệu cho từng tenant.
  \item Mỗi người dùng thuộc một thương hiệu, có một hoặc nhiều vai trò; quyền có thể gán theo vai trò, địa điểm hoặc máy.
  \item Người dùng chỉ truy cập dữ liệu và chức năng trong phạm vi được gán, không truy cập dữ liệu của thương hiệu khác.
  \item Quản trị viên được ủy quyền kích hoạt, khóa, đặt lại tài khoản; phiên đang hoạt động bị thu hồi khi tài khoản bị vô hiệu hóa.
  \item Thao tác nhạy cảm (hoàn tiền, điều chỉnh tồn kho, xịt chẩn đoán, cấu hình lại máy) yêu cầu ủy quyền bổ sung; đăng nhập và thay đổi tài khoản, vai trò, quyền đều được ghi lại.
\end{itemize}
\subsection{Quản lý máy và khe chứa}
\begin{itemize}
  \item Đăng ký mỗi máy bằng số sê-ri và định danh thiết bị duy nhất; liên kết với một tenant và một địa điểm.
  \item Cấu hình số khe trên mỗi máy; mỗi khe có sản phẩm, giá, ngưỡng tồn kho thấp, liều xịt đã hiệu chuẩn và trạng thái sẵn có.
  \item Mỗi khe chỉ có một chai/hộp nước hoa đang hoạt động tại một thời điểm.
  \item Hiển thị trạng thái kết nối, chế độ hoạt động, phiên bản firmware và thời điểm liên lạc cuối; người dùng được ủy quyền bật hoặc tắt máy hoặc từng khe.
  \item Lưu lịch sử lắp đặt, di dời, cấu hình, bảo dưỡng; thay đổi cấu hình quan trọng được ghi vào nhật ký kiểm toán.
\end{itemize}
\subsection{Quản lý sản phẩm nước hoa}

\subsection{Quản lý đơn hàng và thanh toán}
\subsection{Quản lý cấp phát nước hoa}
\subsection{Giám sát thiết bị và dữ liệu IoT}
\subsection{Quản lý tồn kho và tiếp thêm nước hoa}
\subsection{Quản lý cảnh báo và bảo trì}
\subsection{Thống kê, báo cáo và nhật ký hệ thống}

\section{Yêu cầu phi chức năng}

\subsection{Hiệu năng}
\subsection{Tính sẵn sàng}
\subsection{Bảo mật}
\subsection{Độ tin cậy và nhất quán}
\subsection{Khả năng mở rộng}
\subsection{Khả năng sử dụng và bảo trì}
\section{Quy tắc nghiệp vụ}
\begin{enumerate}
\item Mỗi máy thuộc về một thương hiệu tại một thời điểm.
\item Mỗi khe chỉ có một chai/hộp nước hoa đang hoạt động.
\item Đơn hàng lưu giá tại thời điểm tạo.
\item Chỉ backend xác nhận thanh toán thành công.
\item Chỉ tạo lệnh cấp phát sau khi xác nhận thanh toán.
\item Một đơn hàng không có nhiều hơn một lần xịt thành công.
\item Không cấp phát khi máy bảo trì, khe rỗng hoặc có lỗi nghiêm trọng.
\item Kết quả cấp phát không chắc chắn phải được đưa vào kiểm tra.
\item Dữ liệu các thương hiệu phải được phân tách theo quyền.
\end{enumerate}
\section{Biểu đồ Use Case tổng quát}
% TODO: Chèn hình use_case_tong_quat.png khi có.
\section{Đặc tả Use Case tiêu biểu}
\subsection{UC-01: Xem và chọn nước hoa}
\subsection{UC-02: Tạo đơn trải nghiệm}
\subsection{UC-03: Thanh toán đơn hàng}
\subsection{UC-04: Cấp phát một lần xịt}
\subsection{UC-05: Quản lý máy và khe chứa}
\subsection{UC-06: Theo dõi tồn kho và cảnh báo}
\section{Biểu đồ hoạt động}
\subsection{Luồng thanh toán thành công}
\subsection{Luồng thanh toán thất bại hoặc hết hạn}
\subsection{Luồng khe chứa hết nước hoa}
\subsection{Luồng thanh toán thành công nhưng cấp phát lỗi}
\section{Tổng kết chương}
